%=====================================================================
%  ARTIFICIAL INTELLIGENCE --- FALL 2026
%  One-page weekly lecture schedule
%
%  Source of record: parts/G_schedule.tex (weeks, chapters, labs,
%  assessment points) and parts/00_frontmatter.tex (credit hours,
%  prerequisite, suggested weights, set texts).  Change the handout
%  first, then this page.
%
%  Build:  pdflatex AI_Fall2026_Lecture_Schedule.tex
%=====================================================================
\documentclass[10pt,a4paper]{article}
\usepackage[a4paper,left=1.2cm,right=1.2cm,top=0.85cm,bottom=0.6cm]{geometry}
\usepackage[T1]{fontenc}
\usepackage{lmodern}
\usepackage{microtype}
\usepackage[table]{xcolor}
\usepackage{array,ragged2e}
\usepackage{graphicx}
\usepackage{fontawesome5}
\pagestyle{empty}
\pdfinfo{/Title (Artificial Intelligence, Fall 2026: Weekly Lecture Schedule)}
\setlength{\parindent}{0pt}
\setlength{\parskip}{0pt}

%--- palette: the handout's own (aihandout.sty); part colours darkened a
%    little so white chip text stays legible at 6 pt
\definecolor{primary}{HTML}{1A5276}
\definecolor{secondary}{HTML}{2E86C1}
\definecolor{accent}{HTML}{C0392B}
\definecolor{lightbg}{HTML}{EBF5FB}
\definecolor{neutral}{HTML}{707B7C}
\definecolor{sub}{HTML}{465152}
\definecolor{p1}{HTML}{1A5276}
\definecolor{p2}{HTML}{2874A6}
\definecolor{p3}{HTML}{1E8449}
\definecolor{p4}{HTML}{8E44AD}
\definecolor{p5}{HTML}{117A65}
\definecolor{p6}{HTML}{8A6D0A}

%--- building blocks -------------------------------------------------
\newcommand{\chip}[2]{{\setlength{\fboxsep}{1.3pt}\colorbox{#1}{\makebox[1.45em][c]{%
  \fontsize{6.4}{7}\selectfont\bfseries\color{white}#2}}}}
\newcommand{\chipw}[2]{{\setlength{\fboxsep}{1.3pt}\colorbox{#1}{%
  \fontsize{6.4}{7}\selectfont\bfseries\color{white}\,#2\,}}}
% \ch{part}{number}{title}{what it covers}  -- one chapter, own paragraph
\newcommand{\ch}[4]{\chip{p#1}{#2}\,\textbf{#3}\ {\color{neutral}--}\ {\fontsize{7.3}{8.5}\selectfont\color{sub}#4}\par}
% \chn{part}{number}{title}  -- title only, no subtopics
\newcommand{\chn}[3]{\chip{p#1}{#2}\,\textbf{#3}\par}
\newcommand{\ex}[1]{{\bfseries\color{accent}#1}}      % exams and the big deliverables
\newcommand{\ev}[1]{{\bfseries\color{primary}#1}}     % quizzes, assignments, project points
\newcommand{\none}{{\color{neutral}---}}
\newcommand{\hdr}[1]{{\bfseries\color{white}#1}}
\newcommand{\lab}[1]{{\color{sub}#1}}

\newlength{\topicw}   % set inside the document, once \linewidth and \tabcolsep are final
\newcolumntype{K}{>{\centering\arraybackslash}p{0.75cm}}
\newcolumntype{T}{>{\RaggedRight\arraybackslash}p{\topicw}}
\newcolumntype{L}[1]{>{\RaggedRight\arraybackslash}p{#1}}
% \leavevmode: a cell that opens with \color puts a whatsit before its paragraph, and
% \vtop then aligns it by its top edge -- one strut lower than a cell that opens with text
\newcommand{\wk}[4]{\leavevmode{\large\bfseries\color{primary}#1} & \leavevmode #2 &
  \leavevmode #3 & \leavevmode #4 \\ \hline}

\newcommand{\fact}[3]{{\setlength{\fboxsep}{4pt}\colorbox{lightbg}{%
  \parbox[t]{\dimexpr#1-8pt\relax}{\RaggedRight\fontsize{6.2}{7}\selectfont\bfseries\color{neutral}\MakeUppercase{#2}\\[1.5pt]%
  \fontsize{7.8}{9.2}\selectfont\bfseries\color{primary}#3}}}}
\newcommand{\tile}[3]{{\colorbox{lightbg}{\parbox[c][0.9cm][c]{\dimexpr(\linewidth-3pt*\numexpr#1-1\relax)/#1-6pt-0.5pt\relax}{%
  \centering{\large\bfseries\color{primary}#2}\\[-0.5pt]{\fontsize{6.6}{7.4}\selectfont\color{sub}#3}}}}}

% \legend{I}{II}{III}{IV}{V}{VI} -- the six part names, as a 3 x 2 grid of colour chips
\newcommand{\legend}[6]{{\fontsize{6.6}{8}\selectfont\color{sub}%
  \begin{tabular}{@{}p{0.31\linewidth}p{0.31\linewidth}p{0.31\linewidth}@{}}
  \chipw{p1}{I}\ #1 & \chipw{p2}{II}\ #2 & \chipw{p3}{III}\ #3 \\[1pt]
  \chipw{p4}{IV}\ #4 & \chipw{p5}{V}\ #5 & \chipw{p6}{VI}\ #6
  \end{tabular}}}

\begin{document}
\fontsize{8.2}{9.8}\selectfont

%--- header ----------------------------------------------------------
\begin{minipage}[c]{1.5cm}%
  \IfFileExists{slides/assets/iub-logo.png}{\includegraphics[height=1.15cm]{slides/assets/iub-logo.png}}{}%
\end{minipage}\hfill
\begin{minipage}[c]{0.65\linewidth}\centering
  {\footnotesize\bfseries\color{primary}THE ISLAMIA UNIVERSITY OF BAHAWALPUR}\\[1pt]
  {\scriptsize\color{neutral}Department of Information Technology}
\end{minipage}\hfill
\begin{minipage}[c]{1.5cm}\raggedleft
  \IfFileExists{slides/assets/dit-logo.png}{\includegraphics[height=1.15cm]{slides/assets/dit-logo.png}}{}%
\end{minipage}

\vspace{3pt}
{\color{secondary}\rule{\linewidth}{1.4pt}}\\[2pt]
{\LARGE\bfseries\color{primary}ARTIFICIAL INTELLIGENCE}\hfill
{\large\bfseries\color{secondary}Weekly Lecture Schedule\ \textbullet\ Fall 2026}\\[1pt]
{\footnotesize\color{sub}Bachelor of Science in Information Technology\ \textbullet\
 \emph{Search, Knowledge and Reasoning}\ \textbullet\ every example drawn from IT project management}\\[1pt]
{\color{secondary}\rule{\linewidth}{0.6pt}}

\vspace{4pt}
\fact{0.235\linewidth}{Credit hours}{3 (2--3): 2 lecture h + 3 lab h a week}\hfill
\fact{0.235\linewidth}{Duration}{16 teaching weeks; midterm Week 8, final after Week 16}\hfill
\fact{0.235\linewidth}{Prerequisite}{Object Oriented Programming (HEC 2023); Semester 3}\hfill
\fact{0.235\linewidth}{Contents}{22 chapters in 6 parts, each with a runnable Python lab}

\vspace{5pt}
%--- the schedule ----------------------------------------------------
\begingroup   % keep \rowcolors and the table geometry from leaking into the legend
\renewcommand{\arraystretch}{1.3}
\setlength{\tabcolsep}{3.5pt}
\setlength{\topicw}{\dimexpr\linewidth-0.75cm-4.05cm-3.3cm-6\tabcolsep\relax}
\arrayrulecolor{primary!18}
\rowcolors{2}{white}{lightbg}
\noindent\begin{tabular}{@{}K T L{4.05cm} L{3.3cm}@{}}
\rowcolor{primary}
\hdr{\faIcon{calendar-alt}} & \hdr{\faIcon{book-open}\ \ Lecture topics} &
\hdr{\faIcon{flask}\ \ Laboratory (3 h)} & \hdr{\faIcon{clipboard-check}\ \ Assessment} \\
\wk{1}
 {\ch{1}{1}{What Artificial Intelligence Is}{four definitions on two axes, the ``is it AI?'' checklist, branches, five waves, the AI/ML/DS boundary}
  \ch{1}{2}{Intelligent Agents and Task Environments}{percepts and actions, PEAS, the six environment properties, five agent architectures, rationality}}
 {\lab{Is it AI? a five-point audit; a reflex agent on a ticket queue}}{\none}
\wk{2}
 {\ch{1}{3}{Representations and Search in Python}{states as values, the Problem interface, three frontiers, dependency graphs, generators and memoisation}}
 {\lab{The frontier and the graph: the course toolkit}}{\none}
\wk{3}
 {\ch{2}{4}{Problem Formulation and Uninformed Search}{state space against search tree, the explored set, breadth-first, depth-first, iterative deepening, uniform cost}}
 {\lab{Task ordering four ways}}{\ev{Quiz 1} (Ch 1--3)}
\wk{4}
 {\ch{2}{5}{Informed Search and Heuristics}{admissibility and consistency, greedy best-first, $A^*$ and its optimality proof, dominance, relaxation}}
 {\lab{$A^*$ over sprint states, two heuristics}}{\ev{Assignment 1 set}}
\wk{5}
 {\ch{2}{6}{Local Search and Optimisation}{the state-space landscape, hill climbing, random restart, simulated annealing, genetic algorithms}}
 {\lab{Staffing by annealing and by evolution}}{\none}
\wk{6}
 {\ch{3}{7}{Constraint Satisfaction}{variables, domains and constraints, backtracking, minimum remaining values, forward checking, arc consistency}}
 {\lab{Sprint scheduling as a CSP}}{\ev{Assignment 1 due}; \ev{Quiz 2} (Ch 4--6)}
\wk{7}
 {\ch{3}{8}{Adversarial Search and Game Playing}{game trees, minimax, alpha--beta pruning, move ordering, cutting off search, evaluation functions}}
 {\lab{Minimax and alpha--beta on a slot negotiation}}{\none}
\wk{8}
 {\ch{3}{9}{Planning}{the STRIPS representation, progression search, means--ends analysis and the General Problem Solver}}
 {\lab{A release plan from operators}}{\ex{Midterm} (Ch 1--8)}
\wk{9}
 {\ch{4}{10}{Knowledge Representation}{semantic networks, inheritance, frames and defaults, ontologies, what a representation must support}}
 {\lab{A semantic network that answers queries}}{\none}
\wk{10}
 {\ch{4}{11}{Propositional Logic and Inference}{syntax and semantics, entailment, model checking, conjunctive normal form, DPLL, resolution}}
 {\lab{Three ways to decide a release gate}}{\none}
\wk{11}
 {\ch{4}{12}{First-Order Logic and Resolution}{quantifiers, unification with the occurs check, skolemisation, resolution refutation}
  \ch{6}{22}{The Artificial Intelligence Project}{the project brief; choosing and scoping a problem, the plan}}
 {\lab{Unification and a resolution proof}}{\ev{Project set}, proposal due; \ev{Assignment 2 set}}
\wk{12}
 {\ch{4}{13}{Rule-Based Expert Systems}{production rules, forward and backward chaining, conflict resolution, explanation, the knowledge bottleneck}}
 {\lab{A twenty-rule escalation shell that explains itself}}{\ev{Assignment 2 due}}
\wk{13}
 {\ch{5}{14}{Uncertainty and Decision Making}{probability as degree of belief, Bayes' rule, independence, expected utility, the value of information}
  \ch{5}{15}{Bayesian Networks}{structure and conditional independence, enumeration, variable elimination, sampling}}
 {\lab{Go/no-go under uncertainty; diagnosing a defect cause}}{\ev{Quiz 3} (Ch 9--13)}
\wk{14}
 {\ch{5}{16}{Fuzzy Logic}{membership functions, fuzzy operators, Mamdani inference, defuzzification, fuzzy against probabilistic}
  \ch{6}{17}{Learning from Examples}{the three paradigms, the perceptron worked by hand, linear separability and the XOR limit}}
 {\lab{Risk rating by Mamdani inference; the perceptron and its limit}}{\ev{Project experiments frozen}}
\wk{15}
 {\ch{6}{18}{Natural Language Processing}{the symbolic pipeline, context-free grammars, chart parsing, pattern substitution, n-grams}
  \ch{6}{19}{Classic AI Systems and How to Analyse Them}{the General Problem Solver, Eliza, Student, Macsyma; a six-column analysis framework}}
 {\lab{Parsing ticket subjects; four systems in one file}}{\none}
\wk{16}
 {\ch{6}{20}{Recent Trends and Applications}{large language models, tool-using agents, neuro-symbolic systems --- as signposts, not syllabus}
  \ch{6}{21}{Responsible AI and Its Limits}{where bias enters, explainability as an audit artefact, brittleness, computational limits}}
 {\lab{A tool-using agent, offline; auditing the escalation rules}}{\ex{Project report and presentation}}
\end{tabular}
\endgroup

\vspace{5pt}
%--- assessment weights ----------------------------------------------
{\fontsize{6.6}{8}\selectfont\bfseries\color{neutral}\MakeUppercase{Assessment weights}}\\[2pt]
\tile{3}{20\%}{Sessional: quizzes, assignments, project}\hfill
\tile{3}{30\%}{Midterm, Ch 1--8}\hfill
\tile{3}{50\%}{Final, whole course}

\vspace{4pt}
%--- legend and notes ------------------------------------------------
\legend{Foundations of Intelligent Systems}{Problem Solving by Search}{Constraints, Games and Plans}
       {Knowledge, Logic and Reasoning}{Reasoning Under Uncertainty}{Learning, Language and Responsible AI}

\vspace{2pt}
{\fontsize{6.3}{7.6}\selectfont\color{neutral}%
Chapter numbers refer to the Fall 2026 Artificial Intelligence handout (Appendix G, Sixteen-Week Teaching Schedule). Week 8 carries Chapter 9 alongside the midterm on purpose: STRIPS progression search reuses the frontier of Chapters 4 and 5, so the lecture is light in a revision week.
This course is deliberately symbolic --- Chapter 17 is a signpost to \emph{Machine Learning} (Semester 6), not a unit; Appendix F states every boundary, including the one with \emph{Knowledge Representation and Reasoning}.
Weights follow the department's standard split: Sessional 20\%, Mid-term 30\%, Final 50\%. Set texts (HEC): S.\,J. Russell and P.\ Norvig, \emph{Artificial Intelligence: A Modern Approach} (3rd ed., 2010); G.\,F. Luger and W.\,A. Stubblefield, \emph{AI Algorithms, Data Structures, and Idioms} (2009); D.\ Poole and A.\ Mackworth, \emph{Artificial Intelligence: Foundations of Computational Agents} (2nd ed., 2017).\par}

\end{document}
